\section{What the observations establish}
\label{sec:discussion}

\paragraph{Anticipation has an observable temporal structure.}
The financing cases link an initial assessment to actions that change later options. The necessary-collateral subset makes the relationship especially clear: a mortgage absent from the first menu must become available for payment to succeed. Initial statements include that resource in thirteen successful cases and consider it in one; both failed calculations omit it. The successful cases apply a published rule to live holdings, state the resulting resource and proceed to payment.

\paragraph{Local goals and strategic value are different evaluation problems.}
Paying a specified debt supplies a directly checkable endpoint. Bargaining asks whether an exchange is worth making over uncertain future play. The latter depends on rivals' holdings, development, survival and subsequent behavior. G15 records a rival-development objection; G2 combines the same concern with a quantitative error; G19 corrects a pitch and accepts the terms. These observations distinguish expressed consideration, factual accuracy and action consistency. An outcome such as settlement or victory cannot replace a reference value for whether the decision was effective.

Cash has relational value in these interactions. The same transfer can finance the seller's investment, enable the buyer's development and change a third player's prospective inheritance. G7 makes both sides of this calculation explicit: a sale funds Sol's dark blues while the buyer mortgages other deeds to develop orange. G13 instead makes a rival's solvency the object of negotiation. The agents' statements identify these multi-party consequences, even where the records cannot establish that the proposed trade was advantageous. This is the economic content that an agreement rate alone leaves out.

\paragraph{The social behavior is economically situated.}
The conduct spans bluffing, omission, assurance withdrawal and financial-denial appeals. In each case, a statement concerns a concrete resource or opportunity. G13's appeal depends on a third party inheriting the yellows. G18 instead concerns prolonging the last opponent's survival, with that concern expressed alongside the renewed assurance. G7 links a sale to two competing construction programs. These continuations connect the language to the economic stake it addresses.

\paragraph{The interacting system is the unit of observation.}
A provider refusal changes the participant's note and retry path. A conflicting call introduces an offer that another agent may accept. An inaccurate trade summary becomes context for later choices. These conditions propagate through the game, so filtering one participant's rows does not produce an unaffected control group. They are necessary conditions on the interpretation of the behavior, while the executed transactions remain real observations of this cohort.

The broader research questions receive explicit answers in Appendix~\ref{app:questions}. The records document prospective financial assessment, rival-focused bargaining language and conditional communication. G15 shows a blocking policy restated in notes across thirty turns and then applied to three offers; G5 shows a revised inventory followed by payment. Such trajectories provide examples for studying planning, persistence and adaptation.

\section{Related work}
PlanBench tests planning and reasoning about change across domains from automated planning, including the distinction between planning and retrieving familiar knowledge~\cite{valmeekam2023planbench}. Our financing dependency is a published Monopoly rule. The episodes observe its application to a live state, not performance on unfamiliar rules. Refreshed menus and carried notes make this interactive execution rather than generation of an entire plan without feedback. The prospective text is informative because it precedes the release of the needed resource.

Negotiation agents have long been evaluated through both language and outcomes. Lewis et al.~\cite{lewis2017deal} study multi-issue bargaining with private rewards and human negotiation data. He et al.~\cite{he2018decoupling} separate strategic decisions from natural-language generation. Those settings offer explicit utility measures, whereas an offer's value in a complete economic game depends on future interaction. Our distinction between stated assessment and effectiveness follows that established separation.

Guhe and Lascarides~\cite{guhe2014persuasion} use controlled Catan simulations to study when persuasion and trade restrictions improve game outcomes. The strategic idea of withholding trades is therefore not new. Economic models of identity-dependent externalities likewise explain why an asset's buyer can matter to other players~\cite{jehiel1996externalities}. G2 gives a concrete instance of this motive: a non-completing buyer bids to block another player, and the seller quotes that buyer and the prospective completer different prices.

AgentMonopoly compares LLM models and persona instructions through complete games~\cite{wu2026agentmonopoly}. FinBoardBench studies dynamic financial decisions across Cashflow, Acquire and Monopoly, including liquidity management and negotiation~\cite{hu2026finboardbench}. Our contribution is not a new claim that LLMs can play Monopoly. It is the connection between complete negotiation and debt populations, recorded prospective assessments and the particular economic actions they precede.

The Vending-Bench Arena communication study classifies inter-agent messages using simulator state and logged text, with replication across judge families~\cite{li2026arena}. Its population-level communication analysis differs from our purposively selected incidents. We retain contextual transcripts because they make the public claim, private or provider text, recipient response and continuation inspectable. Work on explanation faithfulness further motivates separating recorded explanations from internal computation~\cite{turpin2023unfaithful,lanham2023faithfulness}.

\section{Limitations}
Twenty games span five seeds, one cyclic opponent order and no repeated seed/rotation cells. They cannot separate system, seating, random exposure and opponent effects. Wins are 8/7/3/2 for Sol/Gemini/Opus/Grok and 7/7/3/3 by seat. One seed block is won by Sol from every seat; another is won from seat two by every system. Neither observation establishes a causal decomposition of skill and luck.

Contrasts, cash thresholds and incidents were selected retrospectively. Empirical block reweighting does not correct selection or validate generalization from five blocks. The sixteen necessary-collateral cases form a small retrospective subset with only two failures. There is no evaluated reference policy, counterfactual trade value, independent human annotation or intervention on prompt, memory, communication or response handling. The design does not measure memory benefit, durable learning, calibrated theory of mind or transfer to unfamiliar opponents, and it supplies no independent general-capability axis. Consequently the paper makes no intrinsic model ranking, effectiveness claim for a tactic, population deception estimate or inference about internal reasoning.

\section{Conclusion}
In twenty complete games, LLM agents make prospective statements about finance and rival opportunities, and those statements sometimes align closely with later action. Fourteen of sixteen debts requiring newly released collateral are paid, with initial texts including the resource in thirteen cases and considering it in one. Bargaining shows explicit rival concern alongside opening rejection rates that differ by offer form and change materially with proposer composition. The contextual incidents demonstrate how bluffing, selective presentation, correction and financial denial operate over persistent economic stakes. Together, the records connect stated assessments to economic execution and strategic interaction.
